% =====================================================================
%  APPENDICES
%  Interleave with parent chapters as each becomes final. Most are mechanical
%  (generated from source). Appendix F (synthesis/FPGA) is DROPPED --- FPGA omitted.
%  Use longtable for the big register/descriptor tables.
% =====================================================================

\chapter{CSR register map}
\label{app:csr}

This appendix gives the complete software-visible register map implemented by
\texttt{csr\_registers.sv}, expanding the group overview of \cref{tab:csr-summary}. All
registers are 32 bits wide on a 12-bit byte address; read data is returned one cycle after
the read strobe. Reserved bits are ignored on write and read back as zero.
\Cref{tab:app-csr-ctrl} lists the control, status, and queue-port registers with their
bit-field layouts; \cref{tab:app-csr-timing} lists the timing registers; the Device Address
Table window follows.

\begin{table}[!htb]\centering
  \caption{Control, status, and queue-port registers.}
  \label{tab:app-csr-ctrl}
  \small
  \begin{tabular}{@{}lll>{\raggedright\arraybackslash}p{0.52\textwidth}@{}}
    \toprule
    \textbf{Addr} & \textbf{Register} & \textbf{Access} & \textbf{Bit fields and behaviour} \\
    \midrule
    \texttt{0x004} & \texttt{HC\_CONTROL} & RW &
      [31] \texttt{BUS\_ENABLE} --- enables the protocol FSM;
      [29] \texttt{ABORT} --- host-controller abort request;
      [0] \texttt{IBA\_INCLUDE} --- prefix private I3C transfers with the \texttt{7'h7E}
      broadcast header. Reset \texttt{0x0000\_0000}. \\
    \addlinespace
    \texttt{0x010} & \texttt{RESET\_CONTROL} & W &
      [0] \texttt{SOFT\_RST} --- software-reset request; accepted only while the protocol
      FSM is idle, self-clearing after one cycle, ignored (no-op) while busy. Reads as zero. \\
    \addlinespace
    \texttt{0x014} & \texttt{HC\_STATUS} & R &
      [2] response queue empty; [1] command queue full; [0] protocol FSM idle. \\
    \addlinespace
    \texttt{0x080} & \texttt{CMD\_QUEUE} & W &
      command-descriptor port. Two consecutive 32-bit writes (DWORD~0 then DWORD~1) are
      staged into one 64-bit descriptor and pushed into the Command Queue; the bus stalls
      (\texttt{ready} deasserted) while a completed descriptor awaits queue space. \\
    \addlinespace
    \texttt{0x084} & \texttt{RESP\_PORT} & R &
      response-descriptor pop. Each valid read pops one Response Descriptor
      (\cref{app:descriptors}); reads of an empty queue return zero. \\
    \addlinespace
    \texttt{0x088} & \texttt{PIO\_DATA\_PORT} & RW &
      write: push one 32-bit word into the TX Data Buffer (stalls while a staged word awaits
      space); read: pop one word from the RX Data Buffer (zero when empty). \\
    \addlinespace
    \texttt{0x0B4} & \texttt{QUEUE\_STATUS} & R &
      [7] resp empty; [6] resp full; [5] RX empty; [4] RX full;
      [3] TX empty; [2] TX full; [1] cmd empty; [0] cmd full. \\
    \bottomrule
  \end{tabular}
\end{table}

\begin{table}[!htb]\centering
  \caption{Timing registers. Each holds a 20-bit count of system-clock cycles in bits
    [19:0]; bits [31:20] are reserved. Reset values target \SI{12.5}{\mega\hertz} SDR and
    \SI{400}{\kilo\hertz} \iic{}-FM at the \SI{333}{\mega\hertz} system clock.}
  \label{tab:app-csr-timing}
  \small
  \begin{tabular}{@{}llr@{\hspace{2.5em}}llr@{}}
    \toprule
    \textbf{Addr} & \textbf{Register} & \textbf{Reset} &
    \textbf{Addr} & \textbf{Register} & \textbf{Reset} \\
    \midrule
    \texttt{0x32C} & \texttt{T\_R}          & 4  & \texttt{0x358} & \texttt{I2C\_T\_LOW}    & 534 \\
    \texttt{0x330} & \texttt{T\_F}          & 4  & \texttt{0x35C} & \texttt{T\_HD\_STA}     & 13 \\
    \texttt{0x334} & \texttt{T\_SU\_DAT}    & 1  & \texttt{0x360} & \texttt{I2C\_T\_HD\_STA} & 200 \\
    \texttt{0x338} & \texttt{I2C\_T\_SU\_DAT} & 34 & \texttt{0x368} & \texttt{T\_SU\_STA}   & 7 \\
    \texttt{0x33C} & \texttt{T\_HD\_DAT}    & 0  & \texttt{0x36C} & \texttt{I2C\_T\_SU\_STA} & 200 \\
    \texttt{0x340} & \texttt{T\_HIGH}       & 11 & \texttt{0x370} & \texttt{T\_SU\_STO}     & 7 \\
    \texttt{0x34C} & \texttt{I2C\_T\_HIGH}  & 300 & \texttt{0x374} & \texttt{I2C\_T\_SU\_STO} & 434 \\
    \texttt{0x350} & \texttt{T\_LOW}        & 16 & \texttt{0x37C} & \texttt{T\_BUS\_FREE}   & 13 \\
    \texttt{0x354} & \texttt{T\_LOW\_OD}    & 67 & \texttt{0x380} & \texttt{I2C\_T\_BUF}    & 434 \\
    \bottomrule
  \end{tabular}
\end{table}

\paragraph{Device Address Table window.} The 32 DAT entries occupy word addresses
\texttt{0x400} + $4i$ for $i = 0 \ldots 31$ (\texttt{0x400}--\texttt{0x47C}), reset to zero.
Writes are masked with \texttt{0x807F\_007F}, so only the \texttt{DEVICE} flag and the two
address fields of the entry layout in \cref{app:descriptors} are writable; reads return the
stored entry. The protocol engine reads the same memory through a dedicated hardware port
while executing a command.

\chapter{Command / Response / DAT descriptor formats}
\label{app:descriptors}

This appendix gives the exact bit-field layouts of the descriptors and the address-table
entry summarised in \cref{tab:queue-dat-overview}, taken from \texttt{i3c\_pkg.sv} and
\texttt{controller\_pkg.sv}. Command Descriptors are 64 bits, assembled from two 32-bit
writes to \texttt{CMD\_QUEUE} (DWORD~0 first). The two-bit response-status names are the
\texttt{ERR\_STATUS} encoding of \cref{tab:err-status}.

\paragraph{Command Descriptor DWORD~0 (common fields).} Bits [2:0] carry the
\texttt{ATTR} attribute --- \texttt{Regular} (\texttt{000}), \texttt{Immediate}
(\texttt{001}), \texttt{AddressAssignment} (\texttt{010}), or \texttt{Combo}
(\texttt{011}); bits [6:3] the transaction ID; bits [14:7] the CCC/command byte; bit [15]
the command-present flag (\texttt{CP}); bits [20:16] the DAT device index. The remaining
DWORD-0 bits and the whole of DWORD~1 are attribute-specific, as tabulated below.

\begin{table}[!htb]\centering
  \caption{Attribute-specific fields of the four Command Descriptor variants.}
  \label{tab:app-cmd-desc}
  \small
  \begin{tabular}{@{}l>{\raggedright\arraybackslash}p{0.78\textwidth}@{}}
    \toprule
    \textbf{Attribute} & \textbf{Distinguishing fields} \\
    \midrule
    Regular &
      DWORD0 [21] \texttt{SRE} (report early read end as short-read error), [22] \texttt{DBP}
      (defining byte present), [26:24] \texttt{MODE}, [27] \texttt{RnW}, [28] \texttt{WROC}
      (response on completion), [29] \texttt{TOC} (terminate on completion). DWORD1 [7:0]
      defining byte, [63:48] data length. \\
    \addlinespace
    Immediate &
      DWORD0 [23:21] \texttt{DTT} (number of valid data bytes, 0--4), [26:24] \texttt{MODE},
      [27] \texttt{RnW} (write-only), [28] \texttt{WROC}, [29] \texttt{TOC}. DWORD1 holds up
      to four immediate data bytes (or a defining byte plus three data bytes). \\
    \addlinespace
    AddressAssignment &
      DWORD0 [25:22] \texttt{DEV\_COUNT} (targets to enumerate), [28] \texttt{WROC},
      [29] \texttt{TOC}; DWORD1 reserved. Used to launch ENTDAA. \\
    \addlinespace
    Combo &
      DWORD0 [21] \texttt{FPM}, [22] \texttt{SUB\_16\_OFF}, [24:23] \texttt{DLP},
      [26:24]\footnotemark[1] \texttt{MODE}, [27] \texttt{RnW}, [28] \texttt{WROC},
      [29] \texttt{TOC}. DWORD1 [15:0] offset, [31:16] data length. \\
    \bottomrule
  \end{tabular}
\end{table}
\footnotetext[1]{The Combo layout packs \texttt{DLP}, \texttt{FPM}, and \texttt{MODE} into
the same DWORD-0 region; the field offsets follow \texttt{combo\_trans\_desc\_t} in
\texttt{i3c\_pkg.sv} verbatim.}

\paragraph{Response Descriptor.} One 32-bit word popped from \texttt{RESP\_PORT}:
bits [15:0] the actual data length, bits [23:16] reserved, bits [27:24] the transaction ID
echoed from the command, and bits [31:28] the \texttt{ERR\_STATUS} completion code
(\cref{tab:err-status}).

\paragraph{Device Address Table entry (\texttt{dat\_entry\_t}).} The design uses a 32-bit
entry (halved from the reference's 64-bit entry): bit [31] \texttt{DEVICE} (1 = legacy
\iic{} target), bits [22:16] the I3C dynamic address, bits [6:0] the \iic{} static address,
with bits [30:23] and [15:7] reserved. The write mask \texttt{0x807F\_007F} makes exactly
these three fields writable.

\chapter{Complete FSM state tables}
\label{app:fsm-tables}

This appendix enumerates every state of the seven control FSMs whose diagrams appear in
\cref{ch:architecture-rtl}, taken from the RTL state enumerations and the corresponding
module specifications. The flagship \texttt{flow\_active} table expands
\cref{tab:flow-states}.

\section{\texttt{flow\_active} command engine (\FlowStateCount{} states)}

\begin{longtable}{@{}l>{\raggedright\arraybackslash}p{0.66\textwidth}@{}}
\caption{\texttt{flow\_active} states (enum \texttt{flow\_fsm\_state\_e}, \texttt{flow\_active.sv}).}
\label{tab:app-flow-states}\\
\toprule
\textbf{State (enc.)} & \textbf{Role and principal exits} \\
\midrule
\endfirsthead
\toprule
\textbf{State (enc.)} & \textbf{Role and principal exits} \\
\midrule
\endhead
\bottomrule
\endfoot
\texttt{Idle} (0)              & Wait for the FSM enable; $\to$\texttt{WaitForCmd} on \texttt{i3c\_fsm\_en\_i}. \\
\texttt{WaitForCmd} (1)        & Pop the next Command Descriptor; $\to$\texttt{FetchDAT} when one is available. \\
\texttt{FetchDAT} (2)          & Validate the descriptor and request the DAT lookup; $\to$\texttt{WaitDAT} if supported, $\to$\texttt{WriteResp} (\texttt{NotSupported}) if not. \\
\texttt{WaitDAT} (3)          & One-cycle settle so the captured \texttt{dat\_entry} is current; dispatches to the header, init, or immediate-CCC state by device type and transfer kind. \\
\texttt{I3CBcastHeader} (4)    & Emit \texttt{START + 7'h7E/W + ACK}; branch to \texttt{IssueImmediateCcc}, \texttt{IssueCmd} (ENTDAA), or an I3C init state. \\
\texttt{IssueImmediateCcc} (5) & Execute an immediate ENEC/DISEC CCC; $\to$\texttt{WriteResp} if \texttt{wroc=1}, else $\to$\texttt{Idle}. \\
\texttt{FetchTxData} (6)       & Fetch a TX DWORD; $\to$\texttt{IssueCmd} with data, or $\to$\texttt{WriteResp} when the buffer empties and the STOP completes. \\
\texttt{InitI3CWrite} (7)      & I3C write address phase; $\to$\texttt{FetchTxData} (regular) or \texttt{IssueCmd} (immediate). \\
\texttt{InitI3CRead} (8)       & I3C read address phase; $\to$\texttt{IssueCmd}. \\
\texttt{InitI2CWrite} (9)      & \iic{} write address phase; $\to$\texttt{FetchTxData} or \texttt{IssueCmd}. \\
\texttt{InitI2CRead} (10)      & \iic{} read address phase; $\to$\texttt{IssueCmd}. \\
\texttt{IssueCmd} (11)         & Drive command/data bytes; loops to \texttt{FetchTxData}, continues in-line to \texttt{FetchDAT} on \texttt{toc=0}, self-loops to defer a read abort, or $\to$\texttt{WriteResp}/\texttt{Idle} at completion. \\
\texttt{WriteResp} (12)        & Assert the Response Descriptor and hold until \texttt{resp\_queue\_wready\_i}; $\to$\texttt{Idle}. \\
\end{longtable}

\section{\texttt{scl\_generator} clock/framing engine (14 states)}

\begin{longtable}{@{}l>{\raggedright\arraybackslash}p{0.66\textwidth}@{}}
\caption{\texttt{scl\_generator} states (\texttt{scl\_generator.sv}). \texttt{gen\_idle\_i} forces any state to \texttt{Idle}.}
\label{tab:app-scl-states}\\
\toprule
\textbf{State (enc.)} & \textbf{Role and principal exits} \\
\midrule
\endfirsthead
\toprule
\textbf{State (enc.)} & \textbf{Role and principal exits} \\
\midrule
\endhead
\bottomrule
\endfoot
\texttt{Idle} (0)             & Bus quiescent; $\to$\texttt{GenerateStart} or \texttt{GenerateRstart} on request. \\
\texttt{GenerateStart} (1)    & Set up START; $\to$\texttt{SdaFall} on \texttt{t\_su\_sta}. \\
\texttt{SdaFall} (2)          & Drive SDA low; $\to$\texttt{HoldStart}. \\
\texttt{HoldStart} (3)        & Hold START; $\to$\texttt{DriveLow} on \texttt{t\_hd\_sta}. \\
\texttt{DriveLow} (4)         & SCL low phase; $\to$\texttt{DriveHigh}, \texttt{WaitCmd}, \texttt{GenerateStop}, or \texttt{GenerateRstart} on \texttt{t\_low}. \\
\texttt{DriveHigh} (5)        & SCL high phase; $\to$\texttt{DriveLow}, \texttt{WaitCmd}, \texttt{GenerateStop}, or the Repeated-START states on \texttt{t\_high} (fast path via \texttt{takeover\_i}). \\
\texttt{WaitCmd} (6)          & Hold between bytes; $\to$\texttt{DriveLow}, \texttt{GenerateStop}, or \texttt{GenerateRstart}. \\
\texttt{GenerateRstart} (7)   & Begin Repeated START; $\to$\texttt{SclHighForRstart}. \\
\texttt{SclHighForRstart} (8) & SCL high for Sr; $\to$\texttt{RstartSdaFall} on \texttt{t\_su\_sta}. \\
\texttt{RstartSdaFall} (9)    & SDA fall for Sr; $\to$\texttt{HoldStart}. \\
\texttt{GenerateStop} (10)    & Begin STOP; $\to$\texttt{SclHighForStop}. \\
\texttt{SclHighForStop} (11)  & SCL high for STOP; $\to$\texttt{SdaRise} on \texttt{t\_su\_sto}. \\
\texttt{SdaRise} (12)         & SDA rise (STOP); $\to$\texttt{BusFree}. \\
\texttt{BusFree} (13)         & Enforce bus-free time; $\to$\texttt{Idle} on \texttt{t\_bus\_free} (\texttt{done\_o} pulse). \\
\end{longtable}

\section{ENTDAA subsystem}

\begin{table}[!htb]\centering
  \caption{\texttt{entdaa\_controller} loop manager (7 states) and \texttt{entdaa\_fsm}
    per-Target engine (8 states).}
  \label{tab:app-entdaa-states}
  \small
  \begin{tabular}{@{}l>{\raggedright\arraybackslash}p{0.34\textwidth}
                    l>{\raggedright\arraybackslash}p{0.30\textwidth}@{}}
    \toprule
    \multicolumn{2}{@{}l}{\textbf{\texttt{entdaa\_controller}}} &
    \multicolumn{2}{l}{\textbf{\texttt{entdaa\_fsm}}} \\
    \midrule
    \texttt{Idle} (0)          & await \texttt{daa\_valid\_i}          & \texttt{Idle} (0)         & await \texttt{start\_i} \\
    \texttt{StartLoop} (1)     & more rounds? else done                & \texttt{SendRsvdByte} (1)  & send \texttt{7'h7E/R} \\
    \texttt{RequestRStart} (2) & request Repeated START                & \texttt{ReadRsvdAck} (2)   & ACK $\Rightarrow$ target present \\
    \texttt{WaitRStart} (3)    & await Sr detect                       & \texttt{ReceiveIDBit} (3)  & shift 64 ID bits \\
    \texttt{ReadDAT} (4)       & read next DAT entry                    & \texttt{SendAddr} (4)      & send dyn addr + parity \\
    \texttt{RunEntdaa} (5)     & run child; branch on result           & \texttt{ReadAddrAck} (5)   & ACK $\Rightarrow$ accepted \\
    \texttt{Done} (6)          & $\to$\texttt{Idle}                    & \texttt{WaitStop} (6)      & await STOP \\
                               &                                       & \texttt{Done} (7)          & $\to$\texttt{Idle} \\
    \bottomrule
  \end{tabular}
\end{table}

\section{TX/RX serializers}

\begin{table}[!htb]\centering
  \caption{Bit/byte serializer FSMs: \texttt{bus\_tx} (5 states), \texttt{bus\_tx\_flow}
    (4 states), and \texttt{bus\_rx\_flow} (4 states).}
  \label{tab:app-serdes-states}
  \small
  \begin{tabular}{@{}ll@{\hspace{1.5em}}ll@{\hspace{1.5em}}ll@{}}
    \toprule
    \multicolumn{2}{@{}l}{\textbf{\texttt{bus\_tx}}} &
    \multicolumn{2}{l}{\textbf{\texttt{bus\_tx\_flow}}} &
    \multicolumn{2}{l}{\textbf{\texttt{bus\_rx\_flow}}} \\
    \midrule
    \texttt{Idle}              & idle          & \texttt{Idle}             & idle          & \texttt{Idle}             & idle \\
    \texttt{AwaitClockNegedge} & await SCL     & \texttt{DriveByte}        & 8-bit byte    & \texttt{ReadByte}         & 8-bit byte \\
    \texttt{SetupData}         & setup         & \texttt{DriveBit}         & single bit    & \texttt{ReadBit}          & single bit \\
    \texttt{TransmitData}      & drive on edge & \texttt{NextTaskDecision} & next request  & \texttt{NextTaskDecision} & next request \\
    \texttt{HoldData}          & hold          &                           &               &                           & \\
    \bottomrule
  \end{tabular}
\end{table}

\chapter{CCC subset opcode/frame table}
\label{app:ccc}

This appendix expands the five implemented Common Command Codes of \cref{tab:ccc-supported}
with the on-bus frame each one drives. \texttt{[S]} is START, \texttt{[Sr]} Repeated START,
\texttt{[P]} STOP, \texttt{[7E+W]}/\texttt{[7E+R]} the \texttt{7'h7E} broadcast header with
the direction bit, and \texttt{[a+R/W]} a target's dynamic address. All opcodes match
\texttt{i3c\_pkg::i3c\_ccc\_e}.

\begin{table}[!htb]\centering
  \caption{Implemented CCC opcodes and their bus frames.}
  \label{tab:app-ccc-frames}
  \small
  \begin{tabular}{@{}llc>{\raggedright\arraybackslash}p{0.44\textwidth}@{}}
    \toprule
    \textbf{CCC} & \textbf{Framing} & \textbf{Code} & \textbf{Bus frame} \\
    \midrule
    ENTDAA & Broadcast & \texttt{0x07} &
      \texttt{[S][7E+W][07][Sr][7E+R]} then, per target, 64 ID bits + dynamic address +
      parity + ACK, looping until all assigned, \texttt{[P]}. \\
    \addlinespace
    ENEC   & Broadcast & \texttt{0x00} &
      \texttt{[S][7E+W][00][EVENT][P]} --- enable events on all targets. \\
    DISEC  & Broadcast & \texttt{0x01} &
      \texttt{[S][7E+W][01][EVENT][P]} --- disable events on all targets. \\
    \addlinespace
    ENEC   & Direct    & \texttt{0x80} &
      \texttt{[S][7E+W][80][Sr][a+W][EVENT][P]} --- enable events on one target. \\
    DISEC  & Direct    & \texttt{0x81} &
      \texttt{[S][7E+W][81][Sr][a+W][EVENT][P]} --- disable events on one target. \\
    \bottomrule
  \end{tabular}
\end{table}

The \texttt{EVENT} byte is the CCC-defined event-enable mask; this controller drives it from
the immediate-transfer data byte of the issuing Command Descriptor
(\cref{app:descriptors}).

\chapter{Regression log excerpts}
\label{app:logs}

Each simulation ends with a UVM report summary; a scenario passes when it reports zero
\texttt{UVM\_ERROR} and zero \texttt{UVM\_FATAL}. The listing below is the tail of a
representative run (commit \texttt{047dfdc}, seed~1), showing the four coverage collectors and
the scoreboard reporting cleanly and the zero error/fatal counts that define a pass; every one
of the \RunnableVseqCount{} sequences in \cref{tab:regression-matrix} produced the same
zero-error result.

\begin{lstlisting}[basicstyle=\ttfamily\scriptsize, breaklines=true]
--- UVM Report Summary ---

** Report counts by severity
UVM_INFO  :   25
UVM_WARNING :  0
UVM_ERROR :    0
UVM_FATAL :    0
** Report counts by id
[I3C_CORRELATED_COVERAGE]     1
[I3C_COVERAGE]                1
[I3C_DAA_ARB_COVERAGE]        1
[REG_COVERAGE]                1
[TEST_DONE]                   1
[uvm_test_top.env.m_scoreboard]   16

Simulation complete via $finish(1)
\end{lstlisting}

The merged functional-coverage database is post-processed by the project's own extractor
(\texttt{tools/extract\_coverage.py} over the \texttt{ucd2json} output) into
\texttt{coverage\_report.txt}, whose header records the headline figure and per-coverpoint
detail quoted in \cref{sec:funccov}:

\begin{lstlisting}[basicstyle=\ttfamily\scriptsize, breaklines=true]
I3C Functional Coverage Report
Tests analysed : 80
Total bins     : 1147
Covered bins   : 904
Overall        : 78.8%
\end{lstlisting}

% Appendix F (Synthesis / utilisation / timing) DROPPED --- FPGA omitted.

% Preserve the intentional Appendix-F gap in the printed lettering.
\setcounter{chapter}{6}

\chapter{Build and run instructions}
\label{app:build}

All simulation is driven from the Makefile in \texttt{src/verification/} using Cadence
Xcelium (\texttt{xrun}) with UVM~1.2 (\texttt{CDNS-1.2}). The Xcelium environment must be
sourced once per shell (for example \texttt{source XCELIUM1803.sh}) before any target is
invoked. A single reusable test, \texttt{i3c\_base\_test}, selects the scenario through the
\texttt{+UVM\_TEST\_SEQ} plusarg.

\paragraph{Core targets.}
\begin{itemize}
  \item \texttt{make compile} --- compile and elaborate only.
  \item \texttt{make sim SEQ=\textit{name}} --- run one virtual sequence
    (\cref{app:vseq-inventory}).
  \item \texttt{make sim SEQ=\textit{name} DUMP\_WAVES=1} --- dump an SHM waveform database
    to \texttt{waves/\textit{name}/waves.shm}.
  \item \texttt{make sim SEQ=\textit{name} COV=1} --- collect functional coverage into the
    \texttt{cov\_work} database.
  \item \texttt{VERBOSITY=} and \texttt{SEED=} override the UVM verbosity and simulation
    seed (default seed \texttt{random}; the results of \cref{ch:results} use \texttt{SEED=1}).
\end{itemize}

\paragraph{Regression targets.} Each category target iterates its sequence list; the
aggregate \texttt{regression} target runs all ten in turn, covering the
\RunnableVseqCount{}\,--1 sequences of \cref{app:vseq-inventory} that belong to a target (the
multi-Target arbitration sequence is run on its own).

\begin{table}[!htb]\centering
  \caption{Makefile regression targets and their sequence counts (commit \texttt{047dfdc}).}
  \label{tab:app-make-targets}
  \small
  \begin{tabular}{@{}lcl@{}}
    \toprule
    \textbf{Target} & \textbf{Seqs} & \textbf{Area} \\
    \midrule
    \texttt{csr\_regression}  & 14 & CSR / DAT / register bus \\
    \texttt{fifo\_regression} & 4  & FIFO / queue integrity \\
    \texttt{bus\_regression}  & 11 & PHY / bus / SCL timing \\
    \texttt{sdrw\_regression} & 5  & SDR private write \\
    \texttt{sdrr\_regression} & 6  & SDR private read \\
    \texttt{imm\_regression}  & 4  & immediate transfer \\
    \texttt{i2c\_regression}  & 3  & \iic{}-FM legacy \\
    \texttt{ccc\_regression}  & 5  & CCC frames \\
    \texttt{daa\_regression}  & 6  & ENTDAA / DAA \\
    \texttt{err\_regression}  & 21 & error / status / recovery \\
    \midrule
    \texttt{sdr\_regression}  & \multicolumn{2}{l}{$=$ \texttt{sdrw} $+$ \texttt{sdrr}} \\
    \texttt{regression}       & \multicolumn{2}{l}{runs all ten category targets (79 seqs)} \\
    \bottomrule
  \end{tabular}
\end{table}

\paragraph{Coverage post-processing.} With coverage databases present, \texttt{make cov\_report}
builds the \texttt{ucd2json} tool and runs \texttt{tools/extract\_coverage.py} to produce the
\texttt{coverage\_report.txt} quoted in \cref{sec:funccov,app:logs}. \texttt{make clean}
removes the build, log, waveform, and coverage artifacts.

\chapter{Glossary of I3C and HCI terms}
\label{app:glossary}

This list aligns the thesis terminology with the canonical MIPI I3C vocabulary. Unless noted as
\emph{HCI}, every term is defined in the \textbf{MIPI I3C Basic Specification v1.1.1 (with Errata 01,
2022)}; the clause column gives the defining section. Unlike the front-matter List of
Abbreviations, this appendix supplies definitions and specification pointers rather than only expansions.

{\footnotesize
\begin{longtable}{@{}>{\raggedright\arraybackslash}p{3.0cm}
                    >{\raggedright\arraybackslash}p{1.5cm}
                    >{\raggedright\arraybackslash}p{6.1cm}
                    >{\raggedright\arraybackslash}p{2.5cm}@{}}
\toprule
\textbf{Term} & \textbf{Abbr.} & \textbf{Definition} & \textbf{Spec clause} \\
\midrule
\endfirsthead
\toprule
\textbf{Term} & \textbf{Abbr.} & \textbf{Definition} & \textbf{Spec clause} \\
\midrule
\endhead
\bottomrule
\endfoot
Controller & --- & I3C Device that controls the Bus & \S2.2 \\
Active Controller & --- & The Controller presently in control of the Bus & \S2.2 \\
Target & --- & Device that responds to a Controller & \S2.2 \\
I3C Device & --- & Umbrella term; Controller and Target are Roles & \S2.2 \\
Dynamic Address & --- & Address assigned during Bus initialization & \S2.2 \\
Static Address & --- & Fixed, unchangeable Device Address & \S2.2 \\
Dynamic Address Assignment & DAA & Process of assigning Dynamic Addresses & \S5.1.4 \\
Enter DAA (CCC) & ENTDAA & CCC \texttt{0x07} that starts DAA & \S5.1.9.3.4 \\
Common Command Code & CCC & Standard command issued to Targets & \S2.2, \S5.1.9 \\
Broadcast CCC / Direct CCC & --- & CCC to all Targets / to one Target & \S5.1.9.3 \\
I3C Broadcast Address & --- & \texttt{7'h7E}; addresses all Targets & \S2.2 \\
Defining Byte & --- & Optional byte qualifying a CCC & \S5.1.9.3 \\
Single Data Rate Mode & SDR & Data on one clock edge & \S2.2 \\
High Data Rate Mode & HDR & Higher-speed modes (DDR/TSP/TSL) & \S2.2 \\
Open Drain & OD & High-Z output with active pull-down & \S2.2 \\
Push-Pull & PP & Active pull-up and pull-down driver & \S2.2 \\
START / Repeated START / STOP & Sr & Bus framing conditions & \S2.2 \\
Bus Free Condition & --- & Period after STOP, before START (\texttt{tCAS}/\texttt{tBUF}) & \S5.1.3.2.1 \\
Bus Available Condition & --- & Bus Free sustained \texttt{tAVAL} (Target may START) & \S5.1.3.2.2 \\
Bus Idle Condition & --- & SDA+SCL High for \texttt{tIDLE} (Hot-Join) & \S5.1.3.2.3 \\
Transition Bit & T-Bit & Per-word bit (parity in write, end-of-data in read) replacing ACK in SDR data & \S2.2 \\
Acknowledge / Not-Ack & ACK / NACK & Address-phase acknowledgement & \S2.2 \\
Read/Write bit & RnW & Direction bit in the Address Header & \S5.1.2 \\
Address Header & --- & Address byte + RnW on the Bus & \S5.1.2.2 \\
In-Band Interrupt & IBI & Target emits its address to interrupt the Controller & \S2.2 \\
Hot-Join & --- & Target joins an already-running Bus & \S2.2 \\
Provisioned ID & PID & 48-bit Target identifier used in DAA & \S5.1.4.1.1 \\
Bus Characteristics Register & BCR & 8-bit Target capability register & \S5.1.1.2.1 \\
Device Characteristics Register & DCR & 8-bit Target type register & \S5.1.1.2.2 \\
\midrule
Command / Response Descriptor & --- & HCI: command and completion descriptors & HCI v1.2 / spec 06 \\
Transfer Command / Immediate Data Transfer / Regular Transfer & --- & HCI: command-descriptor transfer types & HCI v1.2 / spec 06 \\
Device Address Table & DAT & HCI: per-Target address/config table & HCI v1.2 / spec 06 \\
Device Characteristics Table & DCT & HCI: DAA result table (unused here) & HCI v1.2 / spec 06 \\
Command / Response Queue & --- & HCI: SW$\to$HW command, HW$\to$SW response queues & HCI v1.2 / spec 06 \\
TX / RX Data Buffer & --- & HCI: SW$\leftrightarrow$HW data queues & HCI v1.2 / spec 06 \\
\midrule
\multicolumn{4}{@{}l@{}}{\itshape Verification (non-MIPI): UVM, TLM, SVA, FSM, CSR --- standard EDA terms, used as-is.} \\
\end{longtable}
}

\chapter{Virtual-sequence inventory}
\label{app:vseq-inventory}

This appendix lists all \RunnableVseqCount{} runnable virtual sequences of
\cref{tab:verif-vseqs}, enumerated from \texttt{i3c\_vseqs/**} at commit \texttt{047dfdc} ---
the baseline of the regression and coverage results in \cref{ch:results}. The count excludes
the four per-category abstract base sequences and the shared ENTDAA stimulus helper, which
are infrastructure rather than runnable tests. The final column gives the owning regression
target from \cref{app:build}; the multi-Target arbitration sequence is runnable on its own but
belongs to no aggregate target.

{\footnotesize
\begin{longtable}{@{}>{\ttfamily}l>{\raggedright\arraybackslash}p{0.44\textwidth}l@{}}
\caption{Runnable virtual-sequence inventory (commit \texttt{047dfdc}).}
\label{tab:app-vseq}\\
\toprule
\normalfont\textbf{Sequence} & \textbf{Objective} & \textbf{Target} \\
\midrule
\endfirsthead
\toprule
\normalfont\textbf{Sequence} & \textbf{Objective} & \textbf{Target} \\
\midrule
\endhead
\bottomrule
\endfoot
\multicolumn{3}{@{}l}{\itshape CSR / DAT / register bus --- \texttt{csr\_regression}} \\
csr\_reset\_defaults\_vseq            & register reset values           & csr \\
csr\_enable\_disable\_vseq            & host-controller enable/disable  & csr \\
csr\_broadcast\_header\_control\_vseq & broadcast-header enable bit     & csr \\
csr\_timing\_rw\_vseq                 & timing-register read/write      & csr \\
csr\_dat\_rw\_all\_entries\_vseq      & DAT read/write, all 32 entries  & csr \\
csr\_dat\_hw\_read\_selection\_vseq   & hardware DAT read-back path     & csr \\
csr\_cmd\_queue\_2dw\_staging\_vseq   & two-DWORD command staging       & csr \\
csr\_cmd\_partial\_then\_other\_write\_vseq & partial descriptor + interleave & csr \\
csr\_rx\_resp\_read\_pop\_vseq        & RX/response port pops           & csr \\
csr\_queue\_status\_flags\_vseq       & queue full/empty status flags   & csr \\
csr\_sw\_reset\_flush\_queues\_vseq   & software reset flushes queues   & csr \\
csr\_sw\_reset\_clears\_cmd\_staging\_vseq & reset clears command staging & csr \\
csr\_unmapped\_addr\_no\_side\_effect\_vseq & unmapped access is inert  & csr \\
csr\_hc\_abort\_control\_vseq         & host-controller abort bit       & csr \\
\midrule
\multicolumn{3}{@{}l}{\itshape FIFO / queue --- \texttt{fifo\_regression}} \\
fifo\_basic\_push\_pop\_order\_vseq   & push/pop ordering               & fifo \\
fifo\_full\_empty\_boundaries\_vseq   & full and empty boundaries       & fifo \\
fifo\_simultaneous\_read\_write\_vseq & concurrent read and write       & fifo \\
fifo\_flush\_during\_activity\_vseq   & flush under active traffic      & fifo \\
\midrule
\multicolumn{3}{@{}l}{\itshape PHY / bus / SCL timing --- \texttt{bus\_regression}} \\
bus\_phy\_reset\_and\_sync\_vseq      & PHY reset and 2FF sync          & bus \\
bus\_start\_stop\_detect\_vseq        & START/STOP detection            & bus \\
bus\_repeated\_start\_detect\_vseq    & Repeated START detection        & bus \\
bus\_monitor\_glitch\_and\_simultaneous\_edge\_filter\_vseq & glitch/edge filtering & bus \\
bus\_monitor\_enable\_gating\_and\_edge\_pulses\_vseq & monitor gating/edge pulses & bus \\
bus\_scl\_clock\_low\_high\_timing\_vseq & SCL low/high timing          & bus \\
bus\_scl\_waitcmd\_stall\_resume\_vseq & WaitCmd stall and resume       & bus \\
bus\_scl\_repeated\_start\_from\_waitcmd\_vseq & Sr from WaitCmd         & bus \\
bus\_tx\_byte\_and\_bit\_order\_vseq  & TX byte/bit order               & bus \\
bus\_rx\_byte\_and\_bit\_order\_vseq  & RX byte/bit order               & bus \\
bus\_tb\_pad\_model\_odpp\_wiring\_vseq & OD/PP pad-model wiring        & bus \\
\midrule
\multicolumn{3}{@{}l}{\itshape SDR private write --- \texttt{sdrw\_regression}} \\
i3c\_write\_vseq                      & baseline private write          & sdrw \\
i3c\_write\_len\_sweep\_vseq          & write length sweep              & sdrw \\
i3c\_write\_toc\_zero\_vseq           & terminate-on-completion = 0     & sdrw \\
i3c\_write\_back\_to\_back\_vseq      & back-to-back writes             & sdrw \\
i3c\_write\_multi\_dat\_idx\_vseq     & multiple DAT indices            & sdrw \\
\midrule
\multicolumn{3}{@{}l}{\itshape SDR private read --- \texttt{sdrr\_regression}} \\
i3c\_read\_vseq                       & baseline private read           & sdrr \\
i3c\_read\_len\_sweep\_vseq           & read length sweep               & sdrr \\
i3c\_read\_target\_more\_than\_requested\_vseq & target over-reads      & sdrr \\
i3c\_read\_toc\_zero\_vseq            & terminate-on-completion = 0     & sdrr \\
i3c\_read\_back\_to\_back\_vseq       & back-to-back reads              & sdrr \\
i3c\_read\_multi\_dat\_idx\_vseq      & multiple DAT indices            & sdrr \\
\midrule
\multicolumn{3}{@{}l}{\itshape Immediate transfer --- \texttt{imm\_regression}} \\
i3c\_imm\_vseq                        & baseline immediate write        & imm \\
i3c\_imm\_dtt\_sweep\_vseq            & data-transfer-type sweep        & imm \\
i3c\_imm\_toc\_vseq                   & immediate terminate-on-completion & imm \\
i3c\_imm\_i2c\_write\_vseq            & \iic{} immediate write          & imm \\
\midrule
\multicolumn{3}{@{}l}{\itshape \iic{} legacy --- \texttt{i2c\_regression}} \\
i2c\_regular\_write\_basic\_vseq      & \iic{} regular write            & i2c \\
i2c\_regular\_read\_basic\_vseq       & \iic{} regular read             & i2c \\
i2c\_len\_sweep\_partial\_rx\_vseq    & partial-RX length sweep         & i2c \\
\midrule
\multicolumn{3}{@{}l}{\itshape CCC frame --- \texttt{ccc\_regression}} \\
i3c\_ccc\_entdaa\_opening\_frame\_vseq & ENTDAA opening frame           & ccc \\
i3c\_ccc\_broadcast\_enec\_vseq       & broadcast ENEC                  & ccc \\
i3c\_ccc\_broadcast\_disec\_vseq      & broadcast DISEC                 & ccc \\
i3c\_ccc\_direct\_enec\_vseq          & direct ENEC                     & ccc \\
i3c\_ccc\_direct\_disec\_vseq         & direct DISEC                    & ccc \\
\midrule
\multicolumn{3}{@{}l}{\itshape ENTDAA / DAA --- \texttt{daa\_regression}} \\
i3c\_daa\_single\_device\_success\_vseq & single-device assignment      & daa \\
i3c\_daa\_no\_device\_vseq            & no responding device            & daa \\
i3c\_daa\_fewer\_devices\_than\_count\_vseq & fewer devices than count  & daa \\
i3c\_daa\_multi\_device\_dat\_loop\_vseq & multi-device DAT loop        & daa \\
i3c\_daa\_address\_rejected\_vseq     & dynamic address rejection       & daa \\
i3c\_daa\_dat\_boundary\_vseq         & DAT boundary index              & daa \\
i3c\_entdaa\_multi\_target\_arbitration\_vseq & wired-AND arbitration   & (standalone) \\
\midrule
\multicolumn{3}{@{}l}{\itshape Error / status / recovery --- \texttt{err\_regression}} \\
i3c\_private\_addr\_nack\_resp\_vseq   & private address NACK            & err \\
i3c\_broadcast\_header\_nack\_resp\_vseq & broadcast-header NACK         & err \\
i3c\_read\_short\_target\_end\_vseq   & early target read end           & err \\
i3c\_write\_tx\_fifo\_underflow\_vseq & TX FIFO underflow               & err \\
i3c\_read\_rx\_fifo\_full\_overflow\_vseq & RX FIFO overflow            & err \\
i3c\_entdaa\_rx\_fifo\_partial\_overflow\_vseq & ENTDAA RX overflow     & err \\
i3c\_resp\_fifo\_full\_backpressure\_vseq & response-FIFO backpressure  & err \\
i3c\_write\_abort\_vseq               & write abort                     & err \\
i3c\_read\_abort\_vseq                & read abort                      & err \\
i3c\_imm\_data\_nack\_i2c\_vseq       & \iic{} immediate data NACK      & err \\
i3c\_imm\_abort\_vseq                 & immediate abort                 & err \\
i2c\_data\_nack\_write\_vseq          & \iic{} write data NACK          & err \\
i2c\_regular\_abort\_vseq             & \iic{} regular abort            & err \\
i3c\_daa\_hc\_abort\_vseq             & abort during DAA                & err \\
i3c\_hc\_abort\_policy\_vseq          & host-controller abort policy    & err \\
i3c\_wroc\_policy\_vseq               & response-on-completion policy   & err \\
i3c\_invalid\_descriptor\_attr\_vseq  & invalid descriptor attribute    & err \\
i3c\_entdaa\_invalid\_descriptor\_resp\_vseq & invalid ENTDAA descriptor & err \\
i3c\_reset\_during\_idle\_vseq        & reset while idle                & err \\
i3c\_reset\_during\_transfer\_phases\_vseq & reset mid-transfer         & err \\
i3c\_sw\_reset\_while\_busy\_policy\_vseq & software-reset-while-busy policy & err \\
\end{longtable}
}
